\section{\nova{}}
\label{sec:method}

\nova{} wraps any probabilistic classifier $f$; we use gradient-boosted trees. It maintains three pieces of
state: a training memory $M$ of labelled flows, a calibration store $Q$ of held-out analyst labels that are never
used for training, and the previously deployed model $f_{\text{prev}}$. \Cref{alg:nova} gives the loop that runs
after every window. Each component addresses one of the four conditions of \cref{sec:intro}:
\begin{itemize}
  \item new classes: open-world scoring with per-class recalibration;
  \item scarce labels: the query and label propagation;
  \item late labels: the memory;
  \item silent degradation: the deployment guard and rollback.
\end{itemize}
\Cref{tab:hyper} lists the hyper-parameters.

\begin{algorithm}[t]
\caption{\nova{} after window $W_w$ has been predicted}
\label{alg:nova}
\begin{algorithmic}[1]
\Require deployed model $f$, previous model $f_{\text{prev}}$, memory $M$, calibration store $Q$
\State $s \gets$ unknown scores of $W_w$ under $f$ \Comment{\cref{sec:unknown}}
\State request labels for $\textsc{Query}(W_w, s, Bn)$; they arrive at window $w{+}\delta$ \Comment{\cref{sec:query}}
\State $L \gets$ analyst labels that arrived now
\If{$f$ was deployed within the last two windows \textbf{and} $\mathrm{F1}_L(f_{\text{prev}}) - \mathrm{F1}_L(f) > \rho$ \textbf{and} $\mathrm{F1}_{\text{ref}}(f) - \mathrm{F1}_{\text{ref}}(f_{\text{prev}}) \le \epsilon$}
   \State roll back to $f_{\text{prev}}$ \Comment{\cref{sec:guard}}
\EndIf
\State $H \gets$ stratified hold-out of $L$ (share $h$); \ $T \gets L\setminus H$
\State $P \gets \textsc{Propagate}(T)$ \Comment{\cref{sec:prop}}
\State $f' \gets \textsc{Train}(M \cup T \cup P)$ with weight $\omega$ on $T \cup P$; calibrate on reference $\cup\; Q$
\If{$\mathrm{F1}_H(f') \ge \mathrm{F1}_H(f)$ \textbf{and} $\mathrm{F1}_{\text{ref}}(f) - \mathrm{F1}_{\text{ref}}(f') \le \epsilon$}
   \State deploy $f'$; $f_{\text{prev}} \gets f$
\EndIf
\State move flows of $H$ to $Q$ if their class has room and $\ge q_{\min}$ flows in $M$, otherwise to $M$
\State $M \gets \textsc{CBRS}(M, T \cup P)$ \Comment{\cref{sec:memory}}
\end{algorithmic}
\end{algorithm}

\subsection{Open-world scoring and recalibration}
\label{sec:unknown}
For a flow $x$ with predicted class $c = \arg\max_j p_j(x)$, \nova{} computes two scores, each normalised so that
its decision threshold is 1:
\begin{equation}
s_{\text{msp}}(x) = \frac{1 - p_c(x)}{1 - \tau^{\text{msp}}_c}, \qquad
s_{\text{knn}}(x) = \frac{d_k\!\left(x, \mathcal{B}_c\right)}{\tau^{\text{knn}}_c},
\end{equation}
where $p_c(x)$ is the maximum softmax probability and $d_k(x,\mathcal{B}_c)$ the Euclidean distance from $x$ to
its $k$-th nearest neighbour in a bank $\mathcal{B}_c$ of up to 3{,}000 training flows of class $c$. The unknown
score is $s(x) = \max\bigl(s_{\text{msp}}(x), s_{\text{knn}}(x)\bigr)$, and the decision is \unk{} if and only if
$s(x) > 1$. The same score ranks flows for the unknown AUROC and drives the query.

\textbf{Per-class thresholds.} A single global threshold is dominated by the majority class: it either rejects
most flows of a newly learned, still sparse class, or accepts everything. \nova{} therefore calibrates one pair
of thresholds per class on held-out flows that the classifier did not train on. $\tau^{\text{msp}}_c$ is the 2\%
quantile of the confidence of held-out flows predicted as $c$; a class with fewer than 20 such flows receives the
most lenient threshold of the other classes. $\tau^{\text{knn}}_c$ is the 99\% quantile of $d_k$ over held-out
flows that belong to $c$ and are predicted as $c$; with fewer than 30 such flows, it falls back to a leave-one-out
estimate on the bank itself.

\textbf{Class-conditional neighbours.} The neighbour search is restricted to the bank of the \emph{predicted}
class. Benign traffic consists of many protocol and application modes, which a single class centroid, as in
CADE~\cite{yang2021cade}, cannot represent: the centroid lies between the modes, and every mode looks distant. A
$k$-NN distance follows each mode. Conditioning on the predicted class also makes the score answer the operational
question, ``is this flow typical of what the model says it is?''.

\textbf{Recalibration.} Thresholds are recomputed for every candidate model on the reference set plus the
calibration store $Q$. Without $Q$, a class that drifts after training keeps thresholds fitted on training-period
traffic. Its later flows are then rejected even after the model has learned the class. This is the mechanism that
pulls the replay baselines below the static model on UNSW-NB15 (\cref{sec:res-main}).

\subsection{Query}
\label{sec:query}
Of the budget $Bn$ labels for a window, a share $\rho_u$ is drawn uniformly at random. This keeps an unbiased
view of benign drift, which a purely novelty-driven query would ignore. The remaining labels are spread over $K$
$k$-means clusters of the window's flows. Cluster $j$ receives at least one label and a quota proportional to
\begin{equation}
q_j \propto \sqrt{|j|}\,\bigl(0.1 + \bar{s}_j\bigr),
\end{equation}
where $|j|$ is the cluster size and $\bar{s}_j$ the mean unknown score of its flows, clipped to $[0,5]$. Inside a
cluster, the flows with the highest unknown score are chosen first.

The square root damps size: a large benign mode still receives labels, but not in proportion to its volume.
The novelty term sends labels to compact clusters of unfamiliar flows, which is how an emerging attack looks in a
window. The constant 0.1 keeps familiar clusters from receiving no labels at all.

\subsection{Label propagation}
\label{sec:prop}
Each analyst label is extended to nearby unlabelled flows of the same window under two agreement conditions. An
unlabelled flow $u$ takes class $c$ when
\begin{enumerate}
  \item its $\min(k_p, n_c)$ nearest labelled flows all have class $c$, where $n_c$ is the number of labels of
        $c$ in the window, and
  \item its distance to the nearest labelled flow is within the radius of $c$, the 90th percentile of
        nearest-same-class distances among the labels of $c$.
\end{enumerate}
At most $r$ propagated labels are kept per analyst label of $c$, closest first, and a class with a single label
does not spread. Unlike self-training~\cite{andresini2021insomnia,zhang2024aoc}, propagation never uses the model's
own predictions and never creates a class; it follows the agreement principle of SOUL~\cite{amalapuram2024soul}.
Hold-out flows are never propagated, so the guard is not evaluated on copies of its own training data.

\subsection{Memory}
\label{sec:memory}
The training memory holds 20{,}000 flows and is a class-balancing reservoir~\cite{chrysakis2020online}. It is
initialised by streaming the training period through the reservoir. When it is full, a new flow of a class that is not the largest evicts
a random flow of the currently largest class, and a flow of the largest class replaces a random member of its
own class with reservoir probability. ``Largest'' compares class size divided by a class weight. Benign traffic
receives weight $\beta$, so it may occupy $\beta$ times the share of any attack class. Benign traffic is
multi-modal, and starving it raises false alarms. The newest labels receive training weight $\omega$, so that a
model update reflects what has just been observed.

\subsection{Deployment guard and rollback}
\label{sec:guard}
Every candidate is trained on the memory plus the new labels, minus a stratified hold-out $H$ of the newest
analyst labels: a share $h$ of each class with at least four labels. The candidate is deployed only if two
conditions hold:
\begin{enumerate}
  \item it is at least as good as the deployed model on $H$;
  \item it loses at most $\epsilon$ macro-F1 on the reference set.
\end{enumerate}
Both comparisons use the open-world credit of \cref{eq:owcredit}, so models that know different classes are
compared on the same task. They ignore classes with fewer than $m$ labels in the compared slice: a class with one
labelled flow makes macro-F1 jump by $1/|\mathcal{C}|$ on a single decision.

After a deployment, the next two label batches are also used to check the new model against its predecessor. If
the predecessor is better by more than $\rho$, \nova{} rolls back. A rollback is itself a deployment and must pass
the same reference check. In our development runs, an unchecked rollback once restored an early model that did not
know most of the attack classes learned since, and undid what the stream had taught.

\subsection{Hyper-parameters and cost}
\Cref{tab:hyper} lists all hyper-parameters. They were set from the design arguments above, checked in two
development runs on CIC-IDS2017, and then frozen for every experiment on every dataset.

Per window, \nova{} performs one $k$-NN scoring pass, one $k$-means run on 4{,}000 flows, and at most one candidate
training on about 20{,}000 flows. The scoring pass uses chunked matrix products against banks of at most 3{,}000
flows per class. Training dominates the cost: \nova{} spends 13--14 minutes of training per 70-window run, 1.2 to 2.2
times as much as CBRS replay (\cref{sec:res-cost}).

\begin{table}[t]
\centering
\caption{\nova{} hyper-parameters (frozen after development on CIC-IDS2017).}
\label{tab:hyper}
\small
\begin{tabular}{@{}llp{0.55\textwidth}@{}}
\toprule
Symbol & Value & Role \\
\midrule
$k$ & 10 & neighbour rank of the $k$-NN score; bank of 3{,}000 flows per class \\
-- & 2\%, 99\% & quantiles for $\tau^{\text{msp}}_c$ and $\tau^{\text{knn}}_c$ \\
$\rho_u$ & 0.2 & uniformly sampled share of the label budget \\
$K$ & 24 & $k$-means clusters for the query \\
$k_p$, $r$ & 5, 5 & propagation agreement size; propagated labels per analyst label \\
$\beta$ & 4 & benign weight in the class-balanced memory (capacity 20{,}000) \\
$\omega$ & 2 & training weight of the newest labels \\
$h$ & 25\% & hold-out share of new analyst labels for the guard \\
$|Q_c|$, $q_{\min}$ & 200, 50 & calibration flows per class; memory flows required before a class feeds $Q$ \\
$\epsilon$ & 0.05 & tolerated reference macro-F1 loss \\
$\rho$ & 0.05 & rollback margin (two-window watch) \\
$m$ & 5 & minimum labels per class in guard comparisons \\
\bottomrule
\end{tabular}
\end{table}
